% !Mode:: "TeX:UTF-8"
\chapter*{致谢}
\addcontentsline{toc}{chapter}{致谢}

时光匆匆，本科阶段的学习与毕业设计工作即将告一段落。回顾本论文从选题、设计、实现到最终成文的过程，其中凝结了许多师长、同学、朋友和家人的帮助与支持。在此，谨向所有关心和帮助过我的人表示诚挚的感谢。

首先，衷心感谢我的指导老师吴际老师。老师在论文选题、研究思路、系统设计和论文写作过程中给予了我耐心而细致的指导。在课题推进过程中，老师帮助我梳理问题边界，在每周的讨论中对我的工作进行点评并给出宝贵的建议；在报告和论文写作方面，老师也对章节结构、实验分析和文字表达提出了许多具体建议，使我能够不断完善文章内容；同时，老师时刻关注着每个时间节点的通知，确保我能够按时完成各个阶段的任务。老师严谨的治学态度和认真负责的工作方式，使我受益良多。

其次，感谢实验室的学长在毕业设计过程中的帮助。在开题阶段，学长向我详细介绍了项目的主要内容、研究目标和整体实现思路，使我能够更快理解课题背景和后续工作的重点。在系统初步完成之后，学长也结合实际功能提出了许多宝贵的改进建议，帮助我发现了一些在开发中容易忽略的问题。与学长的交流让我对项目的实现逻辑和改进方向有了更清晰的认识。

同时，我也要感谢我的家人。感谢你们在我本科四年学习和生活中的理解、支持与包容。无论是在课程学习、项目实践，还是在毕业设计和论文写作阶段，家人的关心都让我能够更加安心地完成眼前的任务。你们的支持是我持续前进的重要依靠。

本科阶段的学习让我逐渐认识到，工程实践不仅需要掌握具体技术，也需要耐心、细致和持续修正问题的能力。本次毕业设计从最初的想法到最终完成，使我对软件系统开发、实验分析和论文写作都有了更具体的体会。未来我仍有许多需要学习和改进的地方，也希望能够带着这段经历中的收获，继续认真面对之后的学习与工作。

\cleardoublepage